\begin{kbvderiv}
\vspace{3pt plus 3pt minus 1pt}
\hspace*{1.8em}ಉದಾಹರಣೆ ; ಗ್‌ + ಅ = ಗ, ಗ್‌ + ಆ = ಗಾ, ಗ್‌ + ಇ = ಗಿ, \\
\hspace*{1.8em}ಗ್‌ + ಈ = ಗೀ ಇತ್ಯಾದಿ. ಇಲ್ಲಿ ಗ್‌ ಅಗ್ರ ಗುಣಿತಾಕ್ಷರಗಳು,
\vspace{3pt plus 3pt minus 1pt}
\end{kbvderiv}

\indent ಈ ರೀತಿಯಾಗಿ ಮೂವತ್ತನಾಲ್ಕು ವ್ಯಂಜನಗಳಿಗೂ ಸ್ವರಗಳನ್ನು ಸೇರಿಸಿ ಗುಣಿತಾಕ್ಷರಗಳನ್ನು ಬರೆಯಬಹುದು. ಬರವಣಿಗೆಯಲ್ಲಿ ಸ್ವರಗಳ ಬದಲು ಸಂಕೇತಗಳನ್ನು ವ್ಯಂಜನಗಳಿಗೆ ಸೇರಿಸಿ ಗುಣಿತಾಕ್ಷರಗಳನ್ನು ಬರೆಯುವುದು ಪದ್ಧತಿ. ಕೆಳಗೆ ಕೊಟ್ಟಿರುವ ಪಟ್ಟಿಯಲ್ಲಿ “ಕ್‌” ವ್ಯಂಜನದ ಗುಣಿತಾಕ್ಷರಗಳನ್ನು ಮಾದರಿಗಾಗಿ ಕೊಟ್ಟಿದೆ.\par

\kbvsechead{ಕಾಗುಣಿತ (ಕ್‌ ವ್ಯಂಜನದ ಗುಣಿತಾಕ್ಷರಗಳು) :}

\vspace{5pt plus 4pt minus 1pt}
\begin{center}
{\renewcommand{\arraystretch}{1.20}\setlength{\tabcolsep}{10pt}
\begin{tabular}{c | c | l}
\toprule
{\kbvsecfont ಸ್ವರ} & {\kbvsecfont ಸ್ವರ ಸಂಕೇತ} & {\kbvsecfont ಗುಣಿತಾಕ್ಷರ} \\
\midrule
ಅ & ಾ (ತಲೆಕಟ್ಟು) & ಕ್‌ + ಅ = ಕ \\
ಆ & ಾ & ಕ್‌ + ಆ = ಕಾ \\
ಇ & ಿ & ಕ್‌ + ಇ = ಕಿ \\
ಈ & ೀ & ಕ್‌ + ಈ = ಕೀ \\
ಉ & ು & ಕ್‌ + ಉ = ಕು \\
ಊ & ೂ & ಕ್‌ + ಊ = ಕೂ \\
ಋ & ೃ & ಕ್‌ + ಋ = ಕೃ \\
ೠ & ೄ & ಕ್‌ + ೠ = ಕೄ \\
ಎ & ೆ & ಕ್‌ + ಎ = ಕೆ \\
ಏ & ೇ & ಕ್‌ + ಏ = ಕೇ \\
ಐ & ೈ & ಕ್‌ + ಐ = ಕೈ \\
ಒ & ೊ & ಕ್‌ + ಒ = ಕೊ \\
ಓ & ೋ & ಕ್‌ + ಓ = ಕೋ \\
ಔ & ೌ & ಕ್‌ + ಔ = ಕೌ \\
ಅಂ & ಂ & ಕ್‌ + ಅಂ = ಕಂ \\
ಅಃ & ಃ & ಕ್‌ + ಅಃ = ಕಃ \\
\bottomrule
\end{tabular}}
\end{center}
\vspace{5pt plus 4pt minus 1pt}
